% =============================================================================
% ch07_oursystem.tex — Chapter 7: NavigatorCrew System | Biomimetic Navigator
% =============================================================================

\selectlanguage{hebrew}
\section{\en{NavigatorCrew} — עיצוב ומימוש הפלטפורמה}
\label{sec:oursystem}

\subsection{סקירת הארכיטקטורה}
פלטפורמת ה-\en{NavigatorCrew} היא מערכת מרובת-סוכנים (\en{Multi-Agent System}) המבוססת על פריימוורק ה-\en{CrewAI}. המערכת נועדה לבצע מחקר, סימולציה וכתיבה של תכנים מדעיים בתחום הניווט האוטונומי בצורה אוטונומית לחלוטין.

\subsection{הסוכנים במערכת}
המערכת מורכבת מחמישה סוכנים מומחים:
\begin{itemize}
    \item \textbf{מנהל המחקר (\en{Research Director}):} אחראי על פירוק נושא המחקר לפרקים והנחיית שאר הסוכנים.
    \item \textbf{חוקר ה-\en{SLAM} והיתוך חיישנים:} מבצע חיפוש אקדמי ב-\en{ArXiv} ו-\en{Serper} ומחלץ אלגוריתמים ונוסחאות.
    \item \textbf{מהנדס הויזואליזציה:} יוצר את הגרפים והדיאגרמות באמצעות קוד \en{Python}.
    \item \textbf{סופר ה-\en{LaTeX}:} מתרגם את ממצאי המחקר לקבצי \en{XeLaTeX} תקניים בעברית ובאנגלית.
    \item \textbf{עורך האיכות:} מבצע ביקורת עמיתים על המאמר, בודק עקביות מתמטית ודיוק בציטוטים.
\end{itemize}

\subsection{הכלים (\en{Tools}) והזיכרון}
הסוכנים מצוידים במערכת כלים המאפשרת להם אינטראקציה עם העולם החיצוני:
\begin{itemize}
    \item \textbf{\en{Search Tools}:} חיפוש ב-\en{ArXiv} ו-\en{Serper.dev}.
    \item \textbf{\en{Code Executor}:} הרצת קוד \en{Python} בסביבה מאובטחת ליצירת גרפים.
    \item \textbf{\en{File Tools}:} קריאה וכתיבה של קבצי טקסט ו-\en{LaTeX}.
\end{itemize}
המערכת משתמשת בזיכרון מבוסס \en{Vector DB} (\en{ChromaDB}) כדי לשמור על עקביות המידע בין הסוכנים לאורך כל שלבי המחקר. \cite{CrewAIDocs}

\subsection{מימוש טכני ודיאגרמת המערכת}
איור \ref{fig:sensor_modalities} (שנמצא בפרק \ref{sec:sensors}) מציג את זרימת המידע, בעוד שאיור \ref{fig:results_summary} בפרק הבא יציג את תוצאות הריצה של המערכת. התשתית ממומשת ב-\en{Python} ורצה על גבי מעבדי \en{NVIDIA Jetson} עבור יישומי זמן-אמת ברחפנים.
